\documentclass{situationist-preprint}

\title{Against Metaethical Neutrality}
\msid{MTE-14-a}
\msfamily{MTE-14}
\msdiscipline{Metaethics}
\mstype{Principal manuscript}
\mskeywords{metaethical neutrality; consequentializing; agent-relative value; fundamentality; reasons first; deontology}

\begin{document}
\maketitle

\begin{abstract}
It is often held that metaethics is neutral between first-order moral theories: that a realist, an expressivist and a constructivist can each be a consequentialist or a deontologist, and that the questions of normative ethics are settled independently of the nature of moral facts. We argue that this is false, and that the consequentializing equivalence of ETH-02 shows exactly why. ETH-02 proves that every deontological theory without dilemmas is extensionally equivalent to a consequentialist theory with an agent-relative value function, and concludes that the dispute between the two is notational. The proof is neutral; the conclusion is not. It holds on a metaethics on which deontic verdicts, or reasons from which both verdicts and values are derived, are the fundamental normative facts. On a metaethics on which agent-neutral value is fundamental and rightness is grounded in it, the agent-relative counterparts are not available, the dispute is substantive, and consequentialism is close to being entailed. Since every complete metaethics takes a stand on the relation between the deontic and the evaluative, every complete metaethics settles whether two first-order theories are one theory or two.
\end{abstract}

\section{The question}

Metaethics asks what moral facts are, how moral language works and how moral knowledge is possible; normative ethics asks which acts are right and which outcomes are good. A widely shared view holds that the two inquiries are independent. On this view, which we call \emph{neutrality}, metaethical theories are compatible with every coherent first-order theory, and whether two first-order theories differ does not depend on which metaethics is true. Neutrality has been defended on the ground that the expressivist can say everything the realist says \citep{Blackburn1993}, and it is presupposed whenever a normative theory is presented as available to realists and anti-realists alike. \citet{Dworkin1996} argued for a stronger thesis, that metaethical claims are themselves first-order moral claims in disguise; \citet{Kramer2009} developed a related view. These positions deny that there is a separate metaethical inquiry to be neutral. We grant that there is such an inquiry and argue that it is not neutral.

The argument is a case study. ETH-02-a proves that every deontological theory without dilemmas is extensionally equivalent to a consequentialist theory with a suitable agent-relative value function, and conversely, and concludes that the dispute between consequentialism and deontology, as a dispute about the right, is notational. We show that the conclusion depends on a metaethical premise, and that the opposite premise reverses it.

\section{Framework}

\begin{definition}\label{d:neutral}
A metaethical theory $M$ is \emph{neutral} with respect to a pair of first-order theories $T_1$ and $T_2$ if $M$ is consistent with each, and whether $T_1$ and $T_2$ are the same theory or rival theories does not depend on whether $M$ is true. Metaethics is \emph{neutral} if every complete metaethical theory is neutral with respect to every pair of first-order theories.
\end{definition}

\begin{definition}\label{d:structure}
A metaethical theory takes a stand on the \emph{order of the normative} if it says which of the following holds: (a) deontic facts (what is permissible, required or forbidden) are fundamental, and evaluative facts about outcomes are grounded in them or defined from them; (b) evaluative facts about outcomes are fundamental, and deontic facts are grounded in them; (c) both are grounded in, or defined from, a common normative primitive, such as reasons; (d) neither is grounded in the other or in a common source.
\end{definition}

\begin{assumption}\label{a:complete}
Every complete metaethical theory takes a stand on the order of the normative.
\end{assumption}

Assumption~\ref{a:complete} is a condition on completeness, not a substantive claim about which stand is correct. A metaethics that said what moral facts are but not how rightness and goodness are related would leave unanswered a question that falls squarely within its subject matter.

\section{What ETH-02 shows}

ETH-02-a establishes, on finite choice structures in which the outcome of an option records the option and the situation, that a deontic theory is extensionally equivalent to some consequentialist theory $C_V$, with $V$ an agent-relative value function, exactly when it permits only available options and generates no dilemmas; that every consequentialist theory is such a theory; and that the two translations are inverse up to equivalence (ETH-02-a, Theorem~3.5 and Corollary~3.6). These are theorems of the mathematics of choice and hold on every metaethics.

The step from equivalence to the conclusion that the dispute is notational is different. ETH-02-a takes two theories to be notational variants with respect to a subject matter if there are translations between them that preserve everything each says about it. The subject matter is fixed as the extension of the right: the data of normative ethics are verdicts about cases. ETH-02-a's reply to the objection from explanatory priority makes the dependence explicit: whether a promise-breaking is wrong because its outcome is worse, or its outcome worse because it is wrong, is a question about the order of grounding between the right and the value of outcomes, and it \enquote{belongs to metaethics}. ETH-02-a's reply to the objection from agent-relative goodness \citep{Schroeder2007} makes a second dependence explicit: the counterparts are consequentialist only if agent-relative rankings of outcomes are rankings by value.

\begin{proposition}\label{p:verdict}
On a metaethics of type (a) or (c), the dispute between consequentialism and deontology is notational in the sense of ETH-02-a.
\end{proposition}

\begin{proof}
On type (a), the fundamental normative facts are deontic, and the value of outcomes is defined from them; an agent-relative value function that ranks the outcomes of permissible options above those of impermissible ones is a legitimate value function, since value just is what the deontic facts make it. Two theories with the same verdicts state the same fundamental facts, and so do not differ. On type (c), say with reasons as the primitive, as MTE-11 holds, the value of an outcome for an agent is defined by the reasons she has to prefer it, and reasons may be agent-relative; so agent-relative value is genuine value and the counterparts are consequentialist. The question whether the right is grounded in the good or the good in the right has a false presupposition, since both are grounded in reasons. A pair of theories with the same verdicts differs, if at all, in which reasons it cites; and a consequentialist counterpart can cite whatever reasons the deontological theory cites, since its value function is defined from them. Either way the residual dispute has no content, and the equivalence makes the dispute notational.
\end{proof}

\begin{proposition}\label{p:value}
On a metaethics of type (b) on which fundamental value is agent-neutral, the dispute is not notational, and a deontological theory that forbids an option with a best outcome is false.
\end{proposition}

\begin{proof}
On such a metaethics, the fundamental evaluative facts rank outcomes in the same way for every agent, and the deontic facts are grounded in them. ETH-02-a, Proposition~4.4, shows that common-sense constraints cannot be represented by an agent-neutral value function on outcomes individuated by the relevant counts; their counterparts need agent-relative values, which on this metaethics are not values but relabelled verdicts. The deontological theory and its counterpart therefore differ in what they say about the ground of rightness, and since the ground of rightness is a fundamental normative fact, the difference is substantive. Moreover, if the deontic facts are grounded in agent-neutral value alone and an option has a best outcome, nothing is available to ground its being forbidden; so a theory that forbids it is false.
\end{proof}

Proposition~\ref{p:value} is close to Moore's own position, on which to assert that an act is right is to assert that it will produce more good than any alternative \citep[ch.~5]{Moore1903}; on that metaethics, consequentialism is not a first-order theory competing with deontology but a consequence of the analysis of \enquote{right}. On type (d), the right and the good are independent, and the correlation that consequentialism asserts between them is a substantive synthetic thesis; the dispute is again not notational, though no side is settled by the metaethics alone.

\begin{theorem}\label{t:main}
Metaethics is not neutral: no complete metaethical theory is neutral with respect to a deontological theory without dilemmas and its consequentialized counterpart.
\end{theorem}

\begin{proof}
By Assumption~\ref{a:complete}, a complete metaethics takes one of the stands (a) to (d). By Proposition~\ref{p:verdict}, on (a) or (c) the two theories are notational variants, and so the same theory. By Proposition~\ref{p:value} and the remark on type (d), on (b) or (d) they are rivals. Since the alternatives are exclusive, whether the two theories are the same depends on which complete metaethics is true, so each complete metaethics fails to be neutral with respect to them. On agent-neutral type (b), in addition, the metaethics is inconsistent with the deontological theory.
\end{proof}

The theorem identifies the precise sense in which ETH-02 illustrates non-neutrality. The equivalence requires a metaethics on which deontic verdicts, or reasons from which verdicts and values are both derived, are fundamental, and value-bearers are not; on the converse metaethics, on which agent-neutral value-bearers are fundamental, the equivalence remains a theorem but no longer shows anything about the dispute, which is substantive and settled in favor of consequentialism. Within the collection, the first condition is met: MTE-11 holds that reasons are primitive and that all other normative notions are defined from them, which is a metaethics of type (c).

\section{Objections}

\begin{objection}[Neutrality about verdicts]
Neutrality, properly stated, is the thesis that metaethics does not settle which acts are right. Theorem~\ref{t:main} concerns whether theories are identical, not which verdicts are correct, and so does not touch it.
\end{objection}

\begin{reply}
Proposition~\ref{p:value} shows that a metaethics of agent-neutral type (b) settles verdicts too: it makes every theory that forbids an option with a best outcome false, and so settles the trolley cases and the cases of constraints against the deontologist. Even the weaker thesis fails. And the identity of theories is not a peripheral matter: whether normative ethics contains a dispute between consequentialism and deontology at all is among the most consequential questions about its shape.
\end{reply}

\begin{objection}[The equivalence is purely formal]
ETH-02's result is a theorem about choice structures; it has no metaethical premises, so it cannot illustrate a dependence on metaethics.
\end{objection}

\begin{reply}
We agree about the theorem and say so above. What depends on metaethics is the step from the theorem to the claim that the dispute is notational, and that step is where ETH-02-a locates its own appeal to metaethics, in its reply to the objection from explanatory priority.
\end{reply}

\begin{objection}[Quasi-realist neutrality]
An expressivist can earn the right to say whatever the realist says, including that the right is grounded in the good. Expressivism therefore takes no stand on the order of the normative and is neutral.
\end{objection}

\begin{reply}
If the expressivist earns the right to say that the right is grounded in the good, then, by her own lights, that is a commitment she has, and it carries the consequences of Proposition~\ref{p:value}. If she declines to say it, she takes another of the stands. A metaethics that can say anything the realist says inherits the realist's need to say one of the things in Definition~\ref{d:structure}. MTE-17 argues, more generally, that once its commitments are explicit, quasi-realism is indistinguishable from realism.
\end{reply}

\section{Conclusion}

ETH-02's equivalence is a theorem of the mathematics of choice, and its conclusion that the dispute between consequentialism and deontology is notational is a consequence of a metaethics on which deontic verdicts, or reasons from which verdicts and values are both derived, are fundamental. On a metaethics on which agent-neutral value is fundamental, the same theorem leaves a substantive dispute, and the metaethics settles it. Since every complete metaethics takes a stand on the order of the normative, every complete metaethics constrains first-order theory, at least by deciding which first-order theories are rivals. Metaethics cannot be neutral between moral theories.

\begin{thebibliography}{99}\small
\bibitem[Blackburn(1993)]{Blackburn1993} Blackburn, S. (1993). \emph{Essays in Quasi-Realism}. Oxford University Press.
\bibitem[Dworkin(1996)]{Dworkin1996} Dworkin, R. (1996). Objectivity and truth: You'd better believe it. \emph{Philosophy \& Public Affairs}, 25(2), 87--139.
\bibitem[Kramer(2009)]{Kramer2009} Kramer, M.~H. (2009). \emph{Moral Realism as a Moral Doctrine}. Wiley-Blackwell.
\bibitem[Moore(1903)]{Moore1903} Moore, G.~E. (1903). \emph{Principia Ethica}. Cambridge University Press.
\bibitem[Schroeder(2007)]{Schroeder2007} Schroeder, M. (2007). Teleology, agent-relative value, and `good'. \emph{Ethics}, 117(2), 265--295.
\end{thebibliography}

\end{document}
